\documentclass[11pt,a4paper]{article}
\usepackage[margin=25mm]{geometry}
\usepackage{amsmath,amssymb,mathtools,enumitem}
\usepackage{xurl}
\usepackage[hidelinks]{hyperref}
\renewcommand{\texttt}[1]{\nolinkurl{#1}}
\setlist[itemize]{nosep,leftmargin=*}
\newcommand{\trans}{\mathsf{T}}
\newcommand{\source}[1]{\par\smallskip\noindent\textit{Source: #1.}\par\smallskip}
\title{EE6221 Robotics\\\large Beginner Notes: Week 1 -- Robot Manipulators and Coordinate Frames}
\author{Living course note}
\date{22 August 2026}
\begin{document}
\maketitle
\noindent\textbf{Scope and sources.} This note covers Week~1 only: robot-manipulator vocabulary, why robots are used, basic specifications, and the coordinate-frame mathematics needed before kinematics modelling. It uses \texttt{resources/lecture-01-manipulator-kinematics.pdf}, pp.~1--52, \texttt{resources/course-outline-part-01.pdf}, and \texttt{resources/week-01-transcript.txt}. The handout controls formulas and terminology.
\tableofcontents

\section{\enspace What is a robot manipulator?}
For this course, a robot is an industrial, reprogrammable, multifunctional manipulator. A manipulator is a chain of rigid \emph{links} joined by actuated \emph{joints}; the final link carries an \emph{end effector}, such as a gripper, hand, or tool. Moving the joints changes the end effector's \emph{pose}: its position and orientation.

The two fundamental joint types are:
\begin{itemize}
\item \textbf{Revolute (R):} the link rotates about an axis; its joint coordinate is an angle.
\item \textbf{Prismatic (P):} the link slides along an axis; its joint coordinate is a displacement.
\end{itemize}
The number of independently controlled joint coordinates is the robot's degrees of freedom (DoF). The first few joints often position the wrist, while wrist joints orient the tool. A robot can be useful with fewer than six DoF, but then it may be unable to realise arbitrary position-and-orientation requests.

\section{\enspace Why and how robots are used}
Robots are particularly useful for repetitive, hazardous, precise, or reconfigurable tasks. In contrast to a hard-automation machine designed for one fixed task, a reprogrammable robot can change task by changing its program and tooling. This flexibility has a cost: it may be less economical than dedicated equipment at very high production volumes, while manual work can be preferable at very low volumes.

Common manipulator structures are named from the major joint types: Cartesian PPP, cylindrical RPP, spherical RRP, SCARA RRP, and articulated RRR. Their structure affects the workspace: the set of tool poses that can be reached. The workspace is not simply a ball around the base. Joint limits, link dimensions, collisions, and orientation requirements restrict it.

Useful specifications include payload (allowable load), reach, stroke, speed, repeatability, accuracy, and resolution. \emph{Accuracy} is closeness to the commanded pose; \emph{repeatability} is how consistently the robot returns to a pose. A robot may be repeatable but inaccurate if a systematic calibration error shifts every placement.
\source{\texttt{resources/lecture-01-manipulator-kinematics.pdf}, pp.~1--22; \texttt{resources/week-01-transcript.txt}.}

\section{\enspace Why coordinate frames are essential}
Saying ``the tool is at $(1,2,3)$'' is incomplete: those numbers only have meaning relative to a coordinate frame. We write $\{A\}$ for a base or reference frame and $\{B\}$ for a frame attached to a robot link. A physical vector is the same geometric object, but its components change when expressed in another frame.

If ${}^{A}\!R_B$ describes frame $\{B\}$ relative to frame $\{A\}$, then
\[
 {}^{A}\!v={}^{A}\!R_B\,{}^{B}\!v.
\]
The columns of ${}^{A}\!R_B$ are the $x_B$, $y_B$, and $z_B$ axes written in frame $\{A\}$. This gives an immediate interpretation: a rotation matrix tells an observer in $\{A\}$ which way the axes attached to $\{B\}$ point.

\section{\enspace Rotations and homogeneous transformations}
A valid three-dimensional rotation matrix satisfies
\[
 R\trans R=I,\qquad R^{-1}=R\trans,\qquad\det R=+1.
\]
For example, a right-handed rotation about the $z$ axis is
\[
 R_z(\theta)=\begin{bmatrix}
 \cos\theta&-\sin\theta&0\\
 \sin\theta&\cos\theta&0\\
 0&0&1
 \end{bmatrix}.
\]
Matrix order matters. Rotations about fixed axes premultiply an existing orientation; rotations about moving (body) axes postmultiply. Always state the frame and physical order instead of relying on a memorised product.

Rotation alone cannot account for different origins. If the origin of $\{B\}$ is at ${}^{A}\!p_B$, a point $P$ satisfies
\[
 {}^{A}\!p_P={}^A\!R_B{}^B\!p_P+{}^A\!p_B.
\]
Homogeneous coordinates combine both operations:
\[
 {}^A\!T_B=\begin{bmatrix}{}^A\!R_B&{}^A\!p_B\\0&1\end{bmatrix},\qquad
 \begin{bmatrix}{}^A\!p_P\\1\end{bmatrix}={}^A\!T_B\begin{bmatrix}{}^B\!p_P\\1\end{bmatrix}.
\]
The inverse is
\[
 ({}^A\!T_B)^{-1}=\begin{bmatrix}({}^A\!R_B)\trans&-({}^A\!R_B)\trans{}^A\!p_B\\0&1\end{bmatrix}.
\]
The translated term must be rotated when inverting; simply negating it is a common mistake.
\source{\texttt{resources/lecture-01-manipulator-kinematics.pdf}, pp.~23--52; \texttt{resources/week-01-transcript.txt}.}

\section{\enspace Week 1 learning checklist}
\begin{itemize}
\item Can I distinguish a link, joint, end effector, pose, and workspace?
\item Can I explain why a reprogrammable robot differs from hard automation?
\item Can I state the source and target frame for every rotation or transformation?
\item Can I explain why the columns of a rotation matrix are axes expressed in another frame?
\item Can I invert a homogeneous transform without forgetting to rotate the translation?
\end{itemize}
\end{document}
